\subsection{Einleitung}
\label{subsec:introduction}

Dieser Bericht dokumentiert die im Rahmen des CUDA-Praktikums entstandene
Implementierung der Canny Edge Detection auf der GPU. Ziel war es, die
einzelnen Stufen des Algorithmus, namentlich Gaussian Blur, Sobel-Filter,
Non-Maximum Suppression und Hysteresis Thresholding, als eigene
CUDA-Kernel umzusetzen und in mehreren Optimierungsstufen (naiv, Shared
Memory, Pinned Memory, sowie stufenspezifische weitere Optimierungen)
gegeneinander zu messen und zu analysieren.

Der Bericht gliedert sich in fünf Abschnitte: Zunächst werden die
Rahmenbedingungen des Projekts sowie die verwendete Messmethodik
beschrieben (Abschnitt~\ref{sec:voraussetzungen-und-projektaufbau}),
gefolgt von einer theoretischen Einordnung des Canny-Algorithmus
(Abschnitt~\ref{sec:algorithmus}). Anschließend werden die konkreten
CUDA-Implementierungen der vier Pipeline-Schritte vorgestellt
(Abschnitt~\ref{sec:implementierung}) und deren Performance anhand von
Laufzeitmessungen sowie NCU-Profiling ausgewertet
(Abschnitt~\ref{sec:ergebnisse-und-analyse}). Den Abschluss bildet ein
gemeinsames Fazit (Abschnitt~\ref{sec:fazit}).